\documentclass[10pt,a4paper]{article}

% Packages essentiels
\usepackage[utf8]{inputenc}
\usepackage[T1]{fontenc}
\usepackage[french]{babel}
\usepackage[margin=1.5cm]{geometry}
\usepackage{amsmath,amssymb}
\usepackage{graphicx}
\usepackage{xcolor}
\usepackage{listings}
\usepackage{tcolorbox}
\tcbuselibrary{breakable}
\usepackage{hyperref}
\usepackage{multicol}

% Chemin des figures
\graphicspath{{./}}

% Configuration listings avec gestion complète des accents
\lstset{
    language=C,
    basicstyle=\ttfamily\scriptsize,
    backgroundcolor=\color{gray!10},
    keywordstyle=\color{blue},
    commentstyle=\color{gray},
    stringstyle=\color{red},
    numbers=left,
    numberstyle=\tiny,
    breaklines=true,
    frame=single,
    tabsize=4,
    showstringspaces=false,
    inputencoding=utf8,
    extendedchars=true,
    literate={é}{{\'e}}1 {è}{{\`e}}1 {à}{{\`a}}1 {ç}{{\c{c}}}1 
             {ê}{{\^e}}1 {î}{{\^i}}1 {ô}{{\^o}}1 {û}{{\^u}}1 {ù}{{\`u}}1
             {É}{{\'E}}1 {È}{{\`E}}1 {À}{{\`A}}1 {Ç}{{\c{C}}}1
             {Ê}{{\^E}}1 {Î}{{\^I}}1 {Ô}{{\^O}}1 {Û}{{\^U}}1 {Ù}{{\`U}}1
             {ë}{{\"e}}1 {ï}{{\"\i}}1 {ö}{{\"o}}1 {ü}{{\"u}}1
             {Ë}{{\"E}}1 {Ï}{{\"I}}1 {Ö}{{\"O}}1 {Ü}{{\"U}}1
             {œ}{{\oe}}1 {Œ}{{\OE}}1 {æ}{{\ae}}1 {Æ}{{\AE}}1
}

% Boxes pédagogiques
\newtcolorbox{definitionbox}[1]{
    colback=blue!5,colframe=blue!75!black,
    fonttitle=\bfseries,title=#1,breakable
}

\newtcolorbox{examplebox}[1]{
    colback=green!5,colframe=green!50!black,
    fonttitle=\bfseries,title=#1,breakable
}

\newtcolorbox{warningbox}[1]{
    colback=red!5,colframe=red!75!black,
    fonttitle=\bfseries,title=#1,breakable
}

\newtcolorbox{techbox}[1]{
    colback=orange!5,colframe=orange!75!black,
    fonttitle=\bfseries,title=#1,breakable
}

\newtcolorbox{infobox}[1]{
    colback=cyan!5,colframe=cyan!75!black,
    fonttitle=\bfseries,title=#1,breakable
}

\title{\textbf{Chapitre 2 : Organisation Mémoire}\\
\large Mots Machine, Adresses, Pointeurs, Tableaux et Chaînes}
\author{Dr. El Hadji Bassirou TOURÉ\\
Département de Mathématiques et Informatique\\
Faculté des Sciences et Techniques\\
Université Cheikh Anta Diop de Dakar}
\date{}

\begin{document}

\maketitle

%==============================================================================
\section{Le Mot Machine et l'Organisation de la Mémoire}
%==============================================================================

\subsection{Définition du Mot Machine}

\begin{definitionbox}{Mot machine (word)}
Un \textbf{mot} (mot machine ou \textit{word}) est l'unité de base manipulée par un microprocesseur. La taille d'un mot s'exprime en bits et sert à classer les architectures : 32 bits ou 64 bits. Cette taille détermine la largeur des registres du processeur, la taille des adresses mémoire, et la quantité de données traitées en un cycle.
\end{definitionbox}

La taille du mot machine a des implications directes sur les capacités du système. Sur une architecture 32 bits, les adresses sont codées sur 32 bits, ce qui limite l'espace d'adressage à $2^{32} = 4$ Go. Sur une architecture 64 bits, l'espace d'adressage théorique est de $2^{64}$ octets, bien que les systèmes actuels n'utilisent que 48 bits d'adresse (256 To).

\subsection{Vue de la Mémoire par un Programme}

Un programme machine perçoit la mémoire comme un très grand tableau linéaire d'octets appelé \textbf{mémoire virtuelle}. Chaque octet possède une adresse unique, et l'ensemble des adresses possibles constitue l'\textbf{espace d'adressage virtuel}.

\begin{center}
\includegraphics[width=0.98\textwidth]{fig01_memory_organization.pdf}
\end{center}

%==============================================================================
\section{Adressage et Représentation des Données}
%==============================================================================

\subsection{Concepts Fondamentaux de l'Adressage}

\begin{definitionbox}{Adresse et Zone Mémoire}
\begin{itemize}
    \item Une \textbf{adresse} est un nombre qui identifie de manière unique un octet en mémoire
    \item Lorsqu'un objet occupe plusieurs octets, son \textbf{adresse} est celle de son premier octet (l'octet à l'adresse la plus basse)
    \item Pour accéder à un objet en mémoire, on a besoin de deux informations : son \textbf{adresse} et sa \textbf{taille}
\end{itemize}
\end{definitionbox}

En C, l'opérateur \texttt{\&} permet d'obtenir l'adresse d'une variable. Cette adresse est toujours celle du premier octet occupé par la variable. Par exemple, si une variable \texttt{int x} est stockée aux adresses 0x1000 à 0x1003, alors \texttt{\&x} vaut \texttt{0x1000}.

\begin{center}
\includegraphics[width=0.85\textwidth]{fig02_address_multioctet.pdf}
\end{center}

\subsection{Taille des Types selon l'Architecture}

La taille des types primitifs en C dépend de l'architecture et du compilateur. Le standard C garantit seulement des tailles minimales et des relations d'ordre : \texttt{sizeof(char) $\leq$ sizeof(short) $\leq$ sizeof(int) $\leq$ sizeof(long)}.

\begin{center}
\includegraphics[width=0.95\textwidth]{fig03_type_sizes.pdf}
\end{center}

\begin{warningbox}{Portabilité}
La différence de taille de \texttt{long} entre 32 et 64 bits est source de bugs. Pour un code portable, utilisez les types de \texttt{<stdint.h>} : \texttt{int32\_t}, \texttt{int64\_t}, \texttt{uint8\_t}, etc.
\end{warningbox}

La taille d'un pointeur correspond \textbf{toujours} à la taille du mot machine : 4 octets sur 32 bits, 8 octets sur 64 bits. C'est logique : un pointeur stocke une adresse, et une adresse occupe exactement un mot machine.

%==============================================================================
\section{Endianness : Ordre des Octets en Mémoire}
%==============================================================================

\subsection{Le Problème de l'Ordre des Octets}

Lorsqu'une donnée occupe plusieurs octets (comme un \texttt{int} de 4 octets), une question se pose : dans quel ordre sont stockés ces octets en mémoire ? Deux conventions existent, appelées \textbf{endianness} (du roman \textit{Les Voyages de Gulliver} de Jonathan Swift).

\begin{definitionbox}{Little-Endian vs Big-Endian}
\begin{itemize}
    \item \textbf{Little-Endian} (``petit-boutiste'') : l'octet de \textbf{poids faible} (LSB - Least Significant Byte) est stocké à l'adresse la plus basse
    \item \textbf{Big-Endian} (``gros-boutiste'') : l'octet de \textbf{poids fort} (MSB - Most Significant Byte) est stocké à l'adresse la plus basse
\end{itemize}
\end{definitionbox}

\subsection{Exemple Détaillé}

Considérons la valeur hexadécimale \texttt{0x01234567} stockée à l'adresse \texttt{0x100}. Cette valeur se décompose en 4 octets : \texttt{01} (MSB), \texttt{23}, \texttt{45}, \texttt{67} (LSB).

\begin{center}
\includegraphics[width=0.95\textwidth]{fig04_endianness.pdf}
\end{center}

\subsection{Quand l'Endianness Importe}

L'endianness devient critique dans trois situations :

\begin{enumerate}
    \item \textbf{Communication réseau} : lorsqu'on envoie des données binaires entre machines de conventions différentes
    \item \textbf{Lecture de fichiers binaires} : un fichier créé sur une machine big-endian sera mal lu sur une machine little-endian
    \item \textbf{Réinterprétation mémoire} : lorsqu'on lit une zone mémoire avec un type différent de celui utilisé pour l'écrire
\end{enumerate}

\begin{lstlisting}[caption={Détection de l'endianness}]
#include <stdio.h>
#include <stdint.h>

int is_little_endian(void) {
    uint32_t x = 1;  /* En mémoire : 01 00 00 00 (LE) ou 00 00 00 01 (BE) */
    uint8_t *p = (uint8_t *)&x;
    return p[0] == 1;  /* Si le premier octet vaut 1, c'est little-endian */
}

int main(void) {
    if (is_little_endian())
        printf("Cette machine est Little-Endian\n");
    else
        printf("Cette machine est Big-Endian\n");
    return 0;
}
\end{lstlisting}

%==============================================================================
\section{Les Pointeurs : Variables Contenant des Adresses}
%==============================================================================

\subsection{Définition et Déclaration}

\begin{definitionbox}{Pointeur}
Un \textbf{pointeur} est une variable dont la valeur est une adresse mémoire. Le type du pointeur indique au compilateur comment interpréter les données à cette adresse et comment effectuer l'arithmétique des pointeurs.
\end{definitionbox}

La déclaration d'un pointeur utilise l'opérateur \texttt{*} :

\begin{lstlisting}
int x = 42;       /* Variable entière */
int *p;           /* Déclaration : p est un pointeur vers un int */
p = &x;           /* p contient maintenant l'adresse de x */

printf("x = %d\n", x);       /* Affiche : x = 42 */
printf("&x = %p\n", &x);     /* Affiche l'adresse de x */
printf("p = %p\n", p);       /* Même adresse que &x */
printf("*p = %d\n", *p);     /* Déréférencement : affiche 42 */
\end{lstlisting}

\subsection{Les Deux Aspects d'un Pointeur}

Comme toute variable, un pointeur possède deux aspects :

\begin{center}
\includegraphics[width=0.95\textwidth]{fig05_pointer_aspects.pdf}
\end{center}

\subsection{Opérateurs sur les Pointeurs}

\begin{techbox}{Opérateurs fondamentaux}
\begin{tabular}{lll}
\texttt{\&x} & Opérateur ``adresse de'' & Retourne l'adresse de la variable x \\
\texttt{*p} & Opérateur de déréférencement & Accède à la valeur à l'adresse p \\
\texttt{p + n} & Addition d'un entier & Avance de n $\times$ sizeof(*p) octets \\
\texttt{p - n} & Soustraction d'un entier & Recule de n $\times$ sizeof(*p) octets \\
\texttt{p - q} & Différence de pointeurs & Nombre d'éléments entre p et q \\
\end{tabular}
\end{techbox}

\subsection{Exemple Pas à Pas : Évolution de la Mémoire}

Suivons l'exécution du code suivant sur une machine 32 bits little-endian :

\begin{lstlisting}
int x = 0x3CD02700;  /* &x = 0x04 */
int y = x + 3;       /* &y = 0x18, y = 0x3CD02703 */
int *z;              /* &z = 0x20 */
z = &y + 3;          /* z = 0x18 + 3*4 = 0x24 */
*z = y;              /* Écrit y à l'adresse 0x24 */
\end{lstlisting}

\begin{center}
\includegraphics[width=0.98\textwidth]{fig06_memory_evolution.pdf}
\end{center}

L'instruction \texttt{z = \&y + 3} illustre l'arithmétique des pointeurs : on ajoute 3 à un pointeur vers \texttt{int}, ce qui ajoute $3 \times 4 = 12$ octets à l'adresse. Ainsi, \texttt{0x18 + 0x0C = 0x24}.

%==============================================================================
\section{Les Tableaux : Zones Mémoire Contiguës}
%==============================================================================

\subsection{Définition et Représentation en Mémoire}

\begin{definitionbox}{Tableau}
Un \textbf{tableau} est une séquence d'emplacements mémoire \textbf{consécutifs} stockant des éléments de même type. Les éléments sont indexés de 0 à $n-1$ où $n$ est la taille du tableau. Un tableau ne connaît pas sa propre taille.
\end{definitionbox}

\begin{lstlisting}
int arr[5] = {10, 20, 30, 40, 50};
/* Si &arr[0] = 0x1000, alors :
   arr[0] est à 0x1000, arr[1] à 0x1004, arr[2] à 0x1008, etc. */
\end{lstlisting}

\begin{center}
\includegraphics[width=0.95\textwidth]{fig07_array_memory.pdf}
\end{center}

\subsection{Absence de Contrôle de Débordement}

\begin{warningbox}{Danger : Pas de vérification des bornes}
Le C ne vérifie \textbf{jamais} si un accès tableau est dans les limites. Un accès hors bornes compile sans erreur et s'exécute, mais corrompt la mémoire ou provoque un crash.
\end{warningbox}

\begin{lstlisting}
int notes[20];
for (int i = 0; i < 40; i++)  /* Boucle jusqu'à 40 au lieu de 20 ! */
    notes[i] = 0;              /* Écrase 80 octets au-delà du tableau */
/* Ce code compile sans warning et peut corrompre la pile */
\end{lstlisting}

Les conséquences d'un débordement dépendent de ce qui se trouve après le tableau en mémoire : autres variables, adresse de retour, etc. C'est une source majeure de vulnérabilités de sécurité (buffer overflow).

\begin{examplebox}{Exemple pas à pas : Tableaux, pointeurs et débordements}
La figure suivante illustre l'évolution de la mémoire lors de manipulations de tableaux et de pointeurs, incluant des accès hors bornes dangereux.
\end{examplebox}

\begin{center}
\includegraphics[width=0.99\textwidth]{fig20_array_overflow_evolution.pdf}
\end{center}

\begin{itemize}
    \item \textbf{Étapes 0-1} : Déclaration de \texttt{int a[6]} et affectations légales
    \item \textbf{Étape 2} : \texttt{a[6] = 0xBAD} écrit \textbf{après} le tableau (buffer overflow)
    \item \textbf{Étape 3} : \texttt{a[-1] = 0xBAD} écrit \textbf{avant} le tableau (aucune erreur !)
    \item \textbf{Étapes 4-5} : Manipulation via pointeur \texttt{p} et arithmétique des pointeurs
\end{itemize}

\subsection{Tableaux à Plusieurs Dimensions}

Un tableau 2D est stocké en mémoire de manière linéaire, ligne par ligne (row-major order).

\begin{lstlisting}
int mat[2][3] = {{1, 2, 3}, {4, 5, 6}};
/* En mémoire : 1, 2, 3, 4, 5, 6 (consécutifs) */
\end{lstlisting}

\begin{center}
\includegraphics[width=0.95\textwidth]{fig08_array_2d.pdf}
\end{center}

%==============================================================================
\section{Relation Fondamentale Pointeurs-Tableaux}
%==============================================================================

\subsection{Conversion Tableau vers Pointeur}

En C, le nom d'un tableau est automatiquement converti en pointeur vers son premier élément dans la plupart des contextes. Cette conversion est appelée \textit{array decay}.

\begin{definitionbox}{Équivalences fondamentales}
\begin{center}
\begin{tabular}{rcl}
\texttt{arr} & $\equiv$ & \texttt{\&arr[0]} \\[3pt]
\texttt{arr + i} & $\equiv$ & \texttt{\&arr[i]} \\[3pt]
\texttt{*(arr + i)} & $\equiv$ & \texttt{arr[i]} \\[3pt]
\texttt{p[i]} & $\equiv$ & \texttt{*(p + i)} \\[3pt]
\end{tabular}
\end{center}
\end{definitionbox}

Ces équivalences expliquent pourquoi la syntaxe \texttt{i[arr]} est valide en C : elle est interprétée comme \texttt{*(i + arr)}, ce qui est identique à \texttt{*(arr + i)} = \texttt{arr[i]}.

\begin{lstlisting}
int arr[5] = {10, 20, 30, 40, 50};
int *p = arr;  /* Équivalent à : int *p = &arr[0]; */

/* Toutes ces expressions sont équivalentes et valent 30 : */
printf("%d\n", arr[2]);     /* Notation tableau */
printf("%d\n", *(arr + 2)); /* Arithmétique pointeur */
printf("%d\n", p[2]);       /* Pointeur avec notation tableau */
printf("%d\n", *(p + 2));   /* Pointeur avec déréférencement */
printf("%d\n", 2[arr]);     /* Syntaxe exotique mais valide ! */
\end{lstlisting}

\subsection{Visualisation de Pointeurs sur un Tableau}

\begin{center}
\includegraphics[width=0.95\textwidth]{fig09_pointers_array.pdf}
\end{center}

\subsection{Différences entre Tableau et Pointeur}

Bien que très similaires dans leur utilisation, tableau et pointeur ne sont pas identiques :

\begin{center}
\includegraphics[width=0.85\textwidth]{fig10_array_vs_pointer.pdf}
\end{center}

%==============================================================================
\section{Classes de Stockage et Durée de Vie des Tableaux}
%==============================================================================

\subsection{Les Trois Classes de Stockage}

En C, un tableau peut être stocké dans trois zones mémoire différentes, chacune avec ses caractéristiques propres.

\begin{center}
\includegraphics[width=0.95\textwidth]{fig11_process_memory.pdf}
\end{center}

\begin{techbox}{Comparaison des classes de stockage}
\begin{tabular}{llll}
\textbf{Classe} & \textbf{Zone} & \textbf{Durée de vie} & \textbf{Taille} \\
\hline
Automatique & Pile (stack) & Durée du bloc \{\} & Fixe à la compilation \\
Statique & Segment data & Toute l'exécution & Fixe à la compilation \\
Dynamique & Tas (heap) & Jusqu'à \texttt{free()} & Variable à l'exécution \\
\end{tabular}
\end{techbox}

\begin{lstlisting}
int global[100];              /* Statique : segment data, existe toujours */

void exemple(int n) {
    int local[10];            /* Automatique : sur la pile, détruit à la fin */
    static int persistant[5]; /* Statique : segment data, conservé entre appels */
    int *dynamique = malloc(n * sizeof(int));  /* Dynamique : sur le tas */
    
    /* ... utilisation ... */
    
    free(dynamique);          /* Libération obligatoire */
}   /* local est détruit ici, persistant survit */
\end{lstlisting}

\subsection{Le Problème du Retour de Tableau}

Une erreur classique consiste à retourner un pointeur vers un tableau local (automatique).

\begin{lstlisting}
/* ERREUR : retourne un pointeur vers une zone mémoire invalide */
int *creer_tableau_FAUX(void) {
    int tab[5] = {1, 2, 3, 4, 5};  /* Alloué sur la pile */
    return tab;  /* Retourne l'adresse de tab */
}   /* tab est détruit ici ! Le pointeur retourné est invalide */

int main(void) {
    int *p = creer_tableau_FAUX();
    printf("%d\n", p[0]);  /* Comportement indéfini ! */
    return 0;
}
\end{lstlisting}

\begin{center}
\includegraphics[width=0.95\textwidth]{fig12_return_local_array.pdf}
\end{center}

\subsection{Solutions pour Retourner un Tableau}

\textbf{Solution 1 : Tableau statique} -- La fonction retourne un pointeur vers un tableau qui persiste.

\begin{lstlisting}
/* OK mais attention : un seul tableau partagé par tous les appels */
int *creer_tableau_static(void) {
    static int tab[5] = {1, 2, 3, 4, 5};  /* Segment data */
    return tab;  /* Valide : tab existe toujours */
}

int main(void) {
    int *p = creer_tableau_static();
    printf("%d\n", p[0]);  /* OK : affiche 1 */
    
    /* ATTENTION : tous les appels partagent le même tableau */
    int *q = creer_tableau_static();
    q[0] = 99;
    printf("%d\n", p[0]);  /* Affiche 99 ! p et q pointent au même endroit */
    return 0;
}
\end{lstlisting}

\textbf{Solution 2 : Allocation dynamique} -- Chaque appel crée un nouveau tableau.

\begin{lstlisting}
/* OK : chaque appel alloue un nouveau tableau */
int *creer_tableau_dynamique(int n) {
    int *tab = malloc(n * sizeof(int));  /* Sur le tas */
    if (tab == NULL) return NULL;        /* Vérifier l'allocation */
    for (int i = 0; i < n; i++)
        tab[i] = i + 1;
    return tab;  /* Valide : le tas n'est pas libéré au retour */
}

int main(void) {
    int *p = creer_tableau_dynamique(5);
    int *q = creer_tableau_dynamique(5);  /* Deux tableaux distincts */
    
    q[0] = 99;
    printf("%d\n", p[0]);  /* Affiche 1 : p et q sont indépendants */
    
    free(p);  /* OBLIGATOIRE : libérer la mémoire */
    free(q);
    return 0;
}
\end{lstlisting}

\textbf{Solution 3 : Passage par paramètre} -- L'appelant fournit le tableau.

\begin{lstlisting}
/* L'appelant gère l'allocation */
void remplir_tableau(int *tab, int n) {
    for (int i = 0; i < n; i++)
        tab[i] = i + 1;
}

int main(void) {
    int local[5];              /* Automatique */
    remplir_tableau(local, 5); /* La fonction remplit le tableau fourni */
    printf("%d\n", local[0]);  /* OK : affiche 1 */
    return 0;
}
\end{lstlisting}

%==============================================================================
\section{Allocation Dynamique : malloc, calloc, realloc, free}
%==============================================================================

\subsection{Principe de l'Allocation Dynamique}

L'allocation dynamique permet de créer des zones mémoire dont la taille est déterminée à l'exécution. Ces zones sont allouées sur le \textbf{tas} (heap) et persistent jusqu'à leur libération explicite.

\begin{techbox}{Fonctions d'allocation (<stdlib.h>)}
\begin{tabular}{ll}
\texttt{malloc(size)} & Alloue \texttt{size} octets (non initialisés) \\
\texttt{calloc(n, size)} & Alloue \texttt{n × size} octets (initialisés à 0) \\
\texttt{realloc(ptr, size)} & Redimensionne la zone pointée par \texttt{ptr} \\
\texttt{free(ptr)} & Libère la zone pointée par \texttt{ptr} \\
\end{tabular}
\end{techbox}

\begin{lstlisting}
#include <stdlib.h>

/* Allocation d'un tableau de n entiers */
int *tab = malloc(n * sizeof(int));  /* n * 4 octets (si int = 4) */

/* TOUJOURS vérifier le retour de malloc */
if (tab == NULL) {
    fprintf(stderr, "Erreur allocation\n");
    return 1;
}

/* Utilisation */
for (int i = 0; i < n; i++)
    tab[i] = i * 10;

/* Libération OBLIGATOIRE */
free(tab);
tab = NULL;  /* Bonne pratique : évite les accès après free */
\end{lstlisting}

\subsection{malloc vs calloc}

\begin{lstlisting}
/* malloc : mémoire NON initialisée (contient des valeurs aléatoires) */
int *p = malloc(5 * sizeof(int));
/* p[0], p[1], ... contiennent des valeurs indéterminées */

/* calloc : mémoire initialisée à ZÉRO */
int *q = calloc(5, sizeof(int));
/* q[0] = q[1] = ... = q[4] = 0 */
\end{lstlisting}

\begin{center}
\includegraphics[width=0.95\textwidth]{fig13_malloc_vs_calloc.pdf}
\end{center}

\subsection{Redimensionnement avec realloc}

\begin{lstlisting}
int *tab = malloc(5 * sizeof(int));
/* ... remplissage de tab[0] à tab[4] ... */

/* Agrandir à 10 éléments */
int *nouveau = realloc(tab, 10 * sizeof(int));
if (nouveau == NULL) {
    /* Échec : tab est toujours valide, mais pas agrandi */
    free(tab);
    return 1;
}
tab = nouveau;  /* Succès : utiliser le nouveau pointeur */
/* tab[0]..tab[4] conservés, tab[5]..tab[9] non initialisés */

free(tab);
\end{lstlisting}

\subsection{Erreurs Classiques}

\begin{lstlisting}
/* 1. Oublier de vérifier malloc */
int *p = malloc(1000000000000);  /* Peut échouer ! */
p[0] = 42;  /* Crash si p == NULL */

/* 2. Oublier free (fuite mémoire) */
void fuite(void) {
    int *p = malloc(100);
    /* ... pas de free(p) ... */
}   /* Mémoire perdue à chaque appel */

/* 3. Double free */
int *p = malloc(10);
free(p);
free(p);  /* Comportement indéfini ! */

/* 4. Accès après free (use-after-free) */
int *p = malloc(10);
free(p);
p[0] = 42;  /* Comportement indéfini ! */
\end{lstlisting}

%==============================================================================
\section{Tableaux à Longueur Variable (VLA)}
%==============================================================================

\subsection{Syntaxe et Utilisation}

Depuis C99, il est possible de déclarer des tableaux dont la taille est une variable (et non une constante).

\begin{lstlisting}
void traiter(int n) {
    int tab[n];  /* VLA : taille déterminée à l'exécution */
    
    for (int i = 0; i < n; i++)
        tab[i] = i * i;
    
    /* tab est automatique : détruit à la fin de la fonction */
}

int main(void) {
    int taille;
    printf("Taille : ");
    scanf("%d", &taille);
    
    int donnees[taille];  /* VLA : alloué sur la pile */
    /* ... */
    return 0;
}
\end{lstlisting}

\subsection{VLA vs Allocation Dynamique}

\begin{center}
\includegraphics[width=0.85\textwidth]{fig14_vla_vs_malloc.pdf}
\end{center}

\begin{warningbox}{Limitations des VLA}
\begin{itemize}
    \item Les VLA sont alloués sur la \textbf{pile}, dont la taille est limitée (typiquement 8 Mo). Un VLA trop grand provoque un \textit{stack overflow}.
    \item Les VLA ne peuvent pas être \texttt{static} ni globaux.
    \item Les VLA sont optionnels depuis C11 (macro \texttt{\_\_STDC\_NO\_VLA\_\_}).
    \item Pour de grands tableaux ou une durée de vie contrôlée, préférer \texttt{malloc}.
\end{itemize}
\end{warningbox}

\begin{lstlisting}
/* Danger : VLA trop grand */
void danger(int n) {
    int enorme[n];  /* Si n = 10000000, stack overflow probable */
}

/* Solution : utiliser malloc pour les grands tableaux */
void safe(int n) {
    int *tab = malloc(n * sizeof(int));
    if (tab == NULL) { /* gérer l'erreur */ }
    /* ... */
    free(tab);
}
\end{lstlisting}

%==============================================================================
\section{Arithmétique des Pointeurs en Détail}
%==============================================================================

\subsection{Règle Fondamentale}

\begin{definitionbox}{Arithmétique dimensionnée}
Toute opération arithmétique sur un pointeur est \textbf{dimensionnée} par la taille du type pointé. Ajouter $n$ à un pointeur de type \texttt{T*} ajoute $n \times \texttt{sizeof(T)}$ octets à l'adresse.
\end{definitionbox}

Cette règle permet de parcourir un tableau naturellement : \texttt{p + 1} pointe toujours vers l'élément suivant, quelle que soit la taille des éléments.

\begin{center}
\includegraphics[width=0.95\textwidth]{fig15_pointer_arithmetic.pdf}
\end{center}

\subsection{Opérations Valides et Exemples}

\begin{lstlisting}
int arr[5] = {10, 20, 30, 40, 50};  /* &arr[0] = 0x1000 */
int *p = arr;       /* p = 0x1000 */
int *q = &arr[3];   /* q = 0x100C */

/* Addition et soustraction d'entiers */
printf("p + 2 = %p\n", (void*)(p + 2));   /* 0x1008 (+8 octets) */
printf("q - 1 = %p\n", (void*)(q - 1));   /* 0x1008 (-4 octets) */

/* Différence de pointeurs : résultat en ÉLÉMENTS, pas en octets */
printf("q - p = %ld\n", q - p);           /* 3 (éléments) */

/* Pour obtenir la différence en octets, cast vers char* */
printf("(char*)q - (char*)p = %ld\n", (char*)q - (char*)p);  /* 12 octets */

/* Comparaisons */
if (p < q) printf("p est avant q en mémoire\n");

/* Incrémentation */
p++;  /* p passe de 0x1000 à 0x1004 */
\end{lstlisting}

\subsection{Réinterprétation Mémoire par Cast}

Un cast de pointeur permet de réinterpréter une zone mémoire avec un type différent. C'est utile pour examiner la représentation interne des données.

\begin{lstlisting}
int x = 0x11223344;
char *cp = (char *)&x;  /* Vue de x comme séquence d'octets */

printf("Octets de x en mémoire:\n");
for (int i = 0; i < 4; i++) {
    printf("  octet %d: 0x%02X\n", i, (unsigned char)cp[i]);
}
/* Sur little-endian : 44 33 22 11 */
/* Sur big-endian    : 11 22 33 44 */
\end{lstlisting}

\begin{center}
\includegraphics[width=0.85\textwidth]{fig19_memory_reinterpretation.pdf}
\end{center}

%==============================================================================
\section{Les Chaînes de Caractères}
%==============================================================================

\subsection{Définition et Représentation}

\begin{definitionbox}{Chaîne de caractères en C}
Une \textbf{chaîne de caractères} est un tableau de \texttt{char} terminé par le caractère nul \texttt{'\textbackslash0'} (code ASCII 0x00). Ce caractère de terminaison permet aux fonctions de déterminer où finit la chaîne.
\end{definitionbox}

\begin{lstlisting}
char str1[] = "Hello";       /* Tableau de 6 chars : H e l l o \0 */
char str2[10] = "Hi";        /* Tableau de 10 chars : H i \0 ? ? ? ? ? ? ? */
char str3[] = {'A', 'B', 0}; /* Équivalent à "AB" */
\end{lstlisting}

\begin{center}
\includegraphics[width=0.95\textwidth]{fig16_string_memory.pdf}
\end{center}

\subsection{Table ASCII Partielle}

Les caractères sont stockés selon leur code ASCII (American Standard Code for Information Interchange).

\begin{center}
\includegraphics[width=0.85\textwidth]{fig17_ascii_table.pdf}
\end{center}

Relations utiles : \texttt{'A'} = 0x41, \texttt{'a'} = 0x61, différence = 0x20. Les chiffres \texttt{'0'} à \texttt{'9'} sont aux codes 0x30 à 0x39.

\subsection{Différence Capitale : Tableau vs Pointeur sur Chaîne}

\begin{lstlisting}
char arr[] = "Hello";    /* Tableau alloué sur la pile, MODIFIABLE */
char *ptr = "Hello";     /* Pointeur vers le segment de données, LECTURE SEULE */

arr[0] = 'J';  /* OK : modifie la copie locale -> "Jello" */
ptr[0] = 'J';  /* ERREUR : segfault ! Le segment données est protégé */

arr = ptr;     /* ERREUR de compilation : arr n'est pas modifiable */
ptr = arr;     /* OK : ptr peut pointer vers une autre adresse */
\end{lstlisting}

\begin{center}
\includegraphics[width=0.98\textwidth]{fig18_array_vs_string_pointer.pdf}
\end{center}

\subsection{Fonctions Standard de Manipulation}

Les fonctions de \texttt{<string.h>} opèrent sur des chaînes terminées par \texttt{'\textbackslash0'}.

\begin{techbox}{Fonctions de <string.h>}
\begin{tabular}{lll}
\texttt{strlen(s)} & Longueur de s & Ne compte pas le \texttt{'\textbackslash0'} \\
\texttt{strcpy(dst, src)} & Copie src dans dst & Danger si dst trop petit \\
\texttt{strncpy(dst, src, n)} & Copie au plus n chars & Plus sûr \\
\texttt{strcat(dst, src)} & Concatène src à dst & Danger si dst trop petit \\
\texttt{strcmp(s1, s2)} & Compare s1 et s2 & 0 si égales, <0 si s1<s2, >0 si s1>s2 \\
\texttt{strchr(s, c)} & Cherche c dans s & Retourne pointeur ou NULL \\
\end{tabular}
\end{techbox}

\begin{lstlisting}
#include <string.h>

char dest[20];
strcpy(dest, "Hello");     /* dest = "Hello" */
strcat(dest, " World");    /* dest = "Hello World" */
printf("Longueur: %zu\n", strlen(dest));  /* 11 */
printf("Taille: %zu\n", sizeof(dest));    /* 20 */

if (strcmp(dest, "Hello World") == 0)
    printf("Chaînes identiques\n");
\end{lstlisting}

%==============================================================================
\section{Visualisation de la Mémoire en C}
%==============================================================================

Le programme suivant permet d'examiner la représentation mémoire des données.

\begin{lstlisting}[caption={Programme complet de visualisation mémoire}]
#include <stdio.h>
#include <stdint.h>

typedef unsigned char *byte_pointer;

/* Affiche len octets à partir de start, 2 caractères par octet */
void show_bytes(byte_pointer start, size_t len) {
    printf("  ");
    for (size_t i = 0; i < len; i++)
        printf("%.2x ", start[i]);
    printf("\n");
}

void show_int(int x) {
    printf("int %d (0x%X):\n", x, x);
    show_bytes((byte_pointer)&x, sizeof(int));
}

void show_float(float x) {
    printf("float %f:\n", x);
    show_bytes((byte_pointer)&x, sizeof(float));
}

void show_pointer(void *x) {
    printf("pointer %p:\n", x);
    show_bytes((byte_pointer)&x, sizeof(void *));
}

int main(void) {
    int val = 0x12345678;
    float f = 1.0f;
    
    show_int(val);
    show_float(f);
    show_pointer(&val);
    
    /* Vérification endianness */
    printf("\nCette machine est ");
    uint32_t test = 1;
    if (*(uint8_t*)&test == 1)
        printf("Little-Endian\n");
    else
        printf("Big-Endian\n");
    
    return 0;
}
\end{lstlisting}

\begin{examplebox}{Sortie typique (64 bits, Little-Endian)}
\begin{verbatim}
int 305419896 (0x12345678):
  78 56 34 12 
float 1.000000:
  00 00 80 3f 
pointer 0x7ffd5c9fa9ac:
  ac a9 9f 5c fd 7f 00 00 

Cette machine est Little-Endian
\end{verbatim}
\end{examplebox}

%==============================================================================
\section{Récapitulatif du Chapitre}
%==============================================================================

\begin{infobox}{Points clés à retenir}
\begin{enumerate}
    \item Le \textbf{mot machine} (32 ou 64 bits) détermine la taille des adresses et des pointeurs
    \item La mémoire est un tableau linéaire d'octets ; chaque octet a une adresse unique
    \item L'\textbf{adresse d'un objet} est l'adresse de son premier octet
    \item \textbf{Endianness} : Little-Endian (Intel) stocke le LSB à l'adresse basse ; Big-Endian fait l'inverse
    \item Un \textbf{tableau} occupe des cases mémoire consécutives ; le C ne vérifie pas les débordements
    \item Le \textbf{nom d'un tableau} est converti en pointeur vers son premier élément
    \item \textbf{Équivalences} : \texttt{arr[i]} $\equiv$ \texttt{*(arr+i)} et \texttt{\&arr[i]} $\equiv$ \texttt{arr+i}
    \item L'\textbf{arithmétique des pointeurs} est dimensionnée par \texttt{sizeof(type pointé)}
    \item Une \textbf{chaîne} est un tableau de \texttt{char} terminé par \texttt{'\textbackslash0'}
    \item \textbf{Tableau vs pointeur sur chaîne} : le tableau est modifiable, le littéral ne l'est pas
\end{enumerate}
\end{infobox}

\begin{techbox}{Formules et syntaxes essentielles}
\begin{tabular}{ll}
Adresse de \texttt{arr[i]} & \texttt{\&arr[0] + i * sizeof(arr[0])} \\
\texttt{ptr + n} & Ajoute \texttt{n * sizeof(*ptr)} octets \\
\texttt{ptr1 - ptr2} & Nombre d'éléments (pas d'octets) \\
\texttt{strlen(s)} & Nombre de caractères avant \texttt{'\textbackslash0'} \\
\texttt{sizeof(arr)} & Taille totale du tableau en octets \\
\texttt{sizeof(ptr)} & Taille du pointeur (4 ou 8) \\
\end{tabular}
\end{techbox}

\end{document}
